% ============================================================================
% Chapter 5: The Credential Verification System
% Study Guide — Blockchain-Based Professional Credential Verification
% ============================================================================

\chapter{The Credential Verification System}
\label{ch:credential-system}

\begin{quote}
\textit{``The preceding chapters established what a blockchain is, how Cardano works, what eUTxO means, and how native tokens and smart contracts operate. This chapter puts all of those pieces together into a single, working application: a credential verification system built on Cardano.''}
\end{quote}

\bigskip

This chapter describes the actual system---every layer, every token, every role, every lifecycle step, and every cost. By the end, you should be able to trace a professional license from its creation on-chain through signing, renewal, revocation, and third-party verification, and explain exactly why each architectural decision was made.

% ============================================================================
\section{System Overview}
\label{sec:ch5-overview}

The credential verification system is organized into three distinct layers, each with a clearly defined responsibility. No single layer can operate in isolation; together they form a complete pipeline from user interaction to immutable on-chain settlement. Understanding this layered separation is essential because it explains why certain operations are fast (local database queries), why others are slow but final (blockchain transactions), and why the middle layer exists at all (to bridge the gap between the two).

\subsection{Layer 1: Cardano Blockchain}

The bottom layer is the Cardano blockchain itself. It provides three critical services that no other layer can replicate: \textbf{token custody} (who holds which tokens is determined by the UTxO set on-chain, not by any off-chain record), \textbf{policy enforcement} (minting policies are evaluated by every node in the network---no single party can bypass them), and \textbf{transaction finality} (once a transaction is confirmed in a block and buried under subsequent blocks, it is irreversible for all practical purposes). The blockchain does not store personal information, does not render dashboards, and does not manage workflow state. It is a settlement layer: slow, expensive relative to a database query, but absolutely authoritative.

Every credential in the system ultimately exists as a Cardano native token. A license is an NFT under a specific minting policy. Signature capacity is a fungible token under the same policy. Validity is a time-bounded token. The blockchain knows nothing about ``licenses'' or ``signatures'' in a semantic sense---it knows only that certain tokens exist at certain addresses under certain policies, and that those policies were satisfied at minting time. All semantic meaning is encoded in metadata and interpreted by the layers above.

\subsection{Layer 2: Off-Chain Application Layer}

The middle layer is a Python application (built on PyCardano v0.19.1) that handles everything the blockchain cannot do efficiently: \textbf{wallet management} (generating HD wallets from mnemonics, deriving keys, encrypting private keys at rest), \textbf{metadata validation} (ensuring CIP-68 metadata conforms to the expected schema before a transaction is built), and \textbf{workflow orchestration} (coordinating multi-step processes like ``mint a license, then mint signature tokens, then mint a validity token, then transfer all three to the licensee''). This layer speaks to the blockchain exclusively through the Blockfrost API, an indexed read/write interface that eliminates the need to run a full Cardano node.

The application layer is where business logic lives. It knows what a ``license'' means, what fields a ``credential'' should have, which wallets are ``authorities'' versus ``licensees,'' and what sequence of transactions constitutes a ``signing ceremony.'' It translates human intent into Cardano transactions and translates on-chain state back into human-readable status.

\subsection{Layer 3: Local SQLite Database}

The top layer is a local SQLite database with nine tables that cache and extend the on-chain state. It exists for three reasons: \textbf{state tracking} (the blockchain is append-only and has no concept of ``current status''---the database maintains derived state like ``this license is active'' by interpreting the chain of transactions), \textbf{caching} (querying the blockchain for every dashboard refresh would be slow and rate-limited; the database provides sub-5ms responses for common queries), and \textbf{dashboard queries} (the web dashboard and CLI tools query the database, not the blockchain, for all read operations). The database is never the source of truth for token ownership---that is always the blockchain. If the database and the chain disagree, the chain wins, and periodic reconciliation corrects the database.

\subsection{Communication: The Blockfrost API}

All communication between the application layer and the Cardano blockchain flows through Blockfrost, a hosted API service that indexes the entire Cardano ledger and exposes it through RESTful endpoints. Blockfrost provides UTxO queries (``what tokens does this address hold?''), transaction submission (``submit this signed transaction to the network''), and chain tip queries (``what is the current slot number?''). The system does not run a full Cardano node. This is a deliberate trade-off: reduced infrastructure cost and complexity in exchange for a dependency on an external service. The risk is mitigated by the existence of multiple Blockfrost alternatives (Koios, Ogmios, self-hosted Blockfrost) and the fact that Blockfrost is a read/write \emph{interface} to the chain, not the chain itself---switching providers does not affect any on-chain state.

\begin{figure}[htbp]
\centering
\begin{tikzpicture}[
  guidebox/.style={draw, rounded corners=4pt, thick, align=center, font=\sffamily\small, fill=white, minimum height=1cm},
  guidearrow/.style={-{Latex[length=2.5mm]}, thick, color=gray!70!black},
  guidelabel/.style={font=\sffamily\footnotesize, text=gray!50!black},
]
  % --- Layer 3: Local Database ---
  \node[guidebox, fill=blue!8, minimum width=3.2cm, text width=3cm] (dash) at (-2.5, 8.5) {Web\\Dashboard};
  \node[guidebox, fill=blue!8, minimum width=3.2cm, text width=3cm] (cli) at (2.5, 8.5) {CLI / Python\\API};

  \node[guidebox, fill=gray!10, minimum width=10cm, minimum height=1.2cm, text width=9.5cm] (db) at (0, 6.8)
    {SQLite Database --- 9 Tables\\[2pt]
    \footnotesize wallets $|$ licenses $|$ signatures $|$ validity $|$ work\_products $|$ policies $|$ dues};
  \node[guidelabel, anchor=east] at (-5.3, 6.8) {\textbf{Layer 3}};

  % --- Layer 2: Application ---
  \node[guidebox, fill=green!8, minimum width=4.5cm, minimum height=1.4cm, text width=4cm] (app1) at (-2.8, 4.8)
    {cardano\_license\\[2pt]\footnotesize Wallet Mgmt $|$ Minting\\Smart Contracts $|$ CIP-68};
  \node[guidebox, fill=green!8, minimum width=4.5cm, minimum height=1.4cm, text width=4cm] (app2) at (2.8, 4.8)
    {tx\_utils\\[2pt]\footnotesize Tx Build $|$ UTxO Select\\Fee Estimate $|$ Submit};
  \node[guidelabel, anchor=east] at (-5.3, 4.8) {\textbf{Layer 2}};

  % --- Blockfrost ---
  \node[guidebox, fill=orange!8, minimum width=10cm, minimum height=1.2cm, text width=9.5cm] (bf) at (0, 2.8)
    {Blockfrost API (Chain Backend)\\[2pt]
    \footnotesize UTxO Queries $|$ Tx Submission $|$ Chain Tip $|$ Asset Info};
  \node[guidelabel, anchor=east] at (-5.3, 2.8) {\textbf{API}};

  % --- Layer 1: Blockchain ---
  \node[guidebox, fill=purple!8, minimum width=10cm, minimum height=1.6cm, text width=9.5cm] (chain) at (0, 0.8)
    {Cardano Blockchain (Layer 1)\\[2pt]
    \footnotesize Token Custody $|$ Policy Enforcement $|$ Transaction Finality\\
    \footnotesize UTxO Ledger $|$ Native Tokens $|$ Plutus V2 Validators};
  \node[guidelabel, anchor=east] at (-5.3, 0.8) {\textbf{Layer 1}};

  % --- Arrows ---
  \draw[guidearrow] (dash.south) -- (db.north -| dash);
  \draw[guidearrow] (cli.south) -- (db.north -| cli);
  \draw[guidearrow] (db.south west) ++(0.5,0) -- (app1.north);
  \draw[guidearrow] (db.south east) ++(-0.5,0) -- (app2.north);
  \draw[guidearrow] (app1.south) -- (bf.north -| app1);
  \draw[guidearrow] (app2.south) -- (bf.north -| app2);
  \draw[guidearrow] (bf.south) -- (chain.north);

  % --- Side annotations ---
  \draw[decorate, decoration={brace, amplitude=6pt, mirror}, thick, gray!50]
    (-5.6, 0) -- (-5.6, 9.1) node[midway, left=8pt, align=right, font=\sffamily\scriptsize, text=gray!50!black]
    {Read: $<$5\,ms (DB)\\Write: 1--3\,s (Tx)\\Finality: $\sim$5\,min};

\end{tikzpicture}
\caption{Three-layer architecture of the credential verification system. Layer~1 (Cardano blockchain) provides settlement finality. Layer~2 (off-chain application) handles business logic and transaction construction. Layer~3 (SQLite) provides fast reads and state caching. All chain communication flows through the Blockfrost API.}
\label{fig:ch5-architecture}
\end{figure}


% ============================================================================
\section{Token Taxonomy}
\label{sec:ch5-tokens}

The system defines four distinct token types, all minted under authority-controlled native script policies. Each token serves a specific purpose in the credential lifecycle, and understanding what each token represents---and who holds it---is fundamental to understanding how the system works. These are not abstract concepts; each token is a real Cardano native asset that occupies a UTxO on-chain, carries a minimum ADA deposit, and can be independently verified by anyone with access to the blockchain.

\subsection{Token 1: License NFT (CIP-68)}

The License NFT is the core credential token. It is \textbf{non-fungible}---each one is unique, representing a specific license held by a specific individual. Unlike simple CIP-25 NFTs (which store metadata in the minting transaction's auxiliary data and are immutable thereafter), the License NFT uses the \textbf{CIP-68 datum metadata standard}, which splits the token into a dual-token pair:

\begin{itemize}[nosep]
  \item \textbf{User Token (label 222):} Held in the licensee's personal wallet. This token is \emph{immutable}---its asset name is fixed at minting time and cannot be changed. It serves as a ``badge'' that the licensee can present to prove they were issued a credential. However, the User Token alone does not prove the credential is \emph{currently valid}---it only proves it was once issued.
  \item \textbf{Reference Token (label 100):} Held at the authority's UTxO (or at a known script address). This token carries an \textbf{inline datum} containing the credential's mutable metadata: license type, jurisdiction, issue date, licensee identity, and---critically---a \texttt{status} field. Because the authority controls the UTxO holding the Reference Token, the authority can update the datum at any time (by spending the UTxO and creating a new one with updated datum) without requiring the licensee's cooperation.
\end{itemize}

The CIP-68 pattern is the architectural foundation of the entire revocation mechanism (detailed in Section~\ref{sec:ch5-revocation}). The key insight is that verification always checks the Reference Token's datum, not the User Token. The User Token in the licensee's wallet is a pointer; the Reference Token at the authority's address is the source of truth.

\subsection{Token 2: Signature Token (Fungible)}

Signature tokens represent \textbf{signing capacity}. They are fungible---one signature token is identical to another under the same policy. When the authority issues a license, it also mints a quantity \texttt{N} of signature tokens and transfers them to the licensee's wallet. Each time the licensee signs a document (or, more precisely, participates in a signing ceremony), one signature token is consumed: it is \textbf{deposited into a validator script address} as part of a transaction. The token is not burned; it is locked at the validator, creating a permanent on-chain record that the signing event occurred.

The quantity \texttt{N} determines how many documents the licensee can sign before needing additional tokens. This creates a natural metering mechanism: a professional engineering board might issue 100 signature tokens per year, corresponding to the expected number of documents a PE would stamp. Running out of signature tokens requires requesting more from the authority, which can impose whatever verification it deems appropriate before minting additional tokens.

\subsection{Token 3: Validity Token (Time-Bounded)}

The validity token represents \textbf{current good standing}. It is a native token with a \textbf{slot-based expiry} encoded in its datum or asset name. The validator script that accepts signature deposits also requires a validity token to accompany each deposit. The validator checks that the validity token's expiry slot is in the future at the time of the transaction. If the validity token has expired, the transaction fails---the licensee cannot sign documents until they obtain a new validity token.

Validity tokens are linked to the dues enforcement contract. When a licensee pays their annual dues (or satisfies whatever renewal criteria the authority requires), the authority mints a new validity token with an expiry slot set to the appropriate future date. This creates a direct, enforceable link between license renewal and signing capacity: if you do not renew, your validity token expires, and you cannot deposit signatures into any validator.

\subsection{Token 4: Work Product (Logical Composite)}

A work product is not a single token but a \textbf{logical composite} tracked in the local SQLite database and coordinated by the \texttt{SignatureCollectionValidator} script. A work product represents a document or project that requires signatures from multiple parties---for example, a set of engineering drawings that need stamps from both the designing engineer and the reviewing engineer.

The validator script address accumulates deposits: each required signer creates a separate UTxO at the script address containing their signature token, their validity token, and a datum identifying the work product (by document hash). When all required signers have deposited, the authority submits a \texttt{Finalize} transaction that consumes all deposit UTxOs atomically. The finalized work product is recorded on-chain as a set of consumed UTxOs under the validator script, with each UTxO containing the signer's public key hash, the document hash, and the deposit slot.

\begin{figure}[htbp]
\centering
\begin{tikzpicture}[
  guidebox/.style={draw, rounded corners=4pt, thick, align=center, font=\sffamily\small, fill=white, minimum height=1cm},
  guidearrow/.style={-{Latex[length=2.5mm]}, thick, color=gray!70!black},
  guidelabel/.style={font=\sffamily\footnotesize, text=gray!50!black},
]
  % === Authority ===
  \node[guidebox, fill=gray!10, minimum width=4cm, minimum height=1.2cm, text width=3.6cm] (auth) at (0, 9.5)
    {\textbf{Authority Wallet}\\[1pt]\footnotesize Minting Policy Owner};

  % === Token 1: License NFT ===
  \node[guidebox, fill=blue!8, minimum width=4.8cm, minimum height=2.8cm, text width=4.4cm] (licnft) at (-4.5, 6)
    {\textbf{Token 1: License NFT}\\(CIP-68 --- Non-Fungible)\\[4pt]
    \footnotesize \textbf{User Token (222)}\\Held by: Licensee\\Immutable badge\\[4pt]
    \textbf{Reference Token (100)}\\Held by: Authority\\Mutable datum:\\status, jurisdiction,\\license type, expiry};

  % === Token 2: Signature Token ===
  \node[guidebox, fill=green!8, minimum width=3.6cm, minimum height=2.2cm, text width=3.2cm] (sigtok) at (1, 6)
    {\textbf{Token 2: Signature}\\(Fungible)\\[4pt]
    \footnotesize Held by: Licensee\\Qty: N per licensee\\Consumed: 1 per\\signing operation\\Deposited to validator};

  % === Token 3: Validity Token ===
  \node[guidebox, fill=orange!8, minimum width=3.6cm, minimum height=2.2cm, text width=3.2cm] (valtok) at (5.5, 6)
    {\textbf{Token 3: Validity}\\(Time-Bounded)\\[4pt]
    \footnotesize Held by: Licensee\\Slot-based expiry\\Must accompany\\signature deposits\\Renewed via dues};

  % === Licensee Wallet ===
  \node[guidebox, fill=yellow!8, minimum width=4cm, minimum height=1cm, text width=3.6cm] (licensee) at (3.2, 2.8)
    {\textbf{Licensee Wallet}\\[1pt]\footnotesize Holds: User Token + Sig Tokens + Val Token};

  % === Token 4: Work Product / Validator ===
  \node[guidebox, fill=purple!8, minimum width=10cm, minimum height=1.8cm, text width=9.5cm] (wp) at (0.5, 0.5)
    {\textbf{Token 4: Work Product} (Logical Composite)\\[3pt]
    \footnotesize \textbf{SignatureCollectionValidator} script address\\
    Collects sig + val tokens from each required signer\\
    Redeemers: Collect (signer) $|$ Finalize (authority) $|$ Reclaim (cancel)\\
    Each signer $\to$ independent UTxO (avoids eUTxO contention)};

  % === Arrows: Authority mints ===
  \draw[guidearrow] (auth.south) -- ++(0,-0.5) -| (licnft.north);
  \draw[guidearrow] (auth.south) -- (sigtok.north);
  \draw[guidearrow] (auth.south) -- ++(0,-0.5) -| (valtok.north);

  % === Arrows: Tokens to licensee ===
  \draw[guidearrow, dashed, blue!60!black] (licnft.south) -- ++(0,-0.3)
    node[below, font=\sffamily\tiny, text=blue!60!black] {User Token (222)} -| (licensee.west);
  \draw[guidearrow, dashed, green!60!black] (sigtok.south) --
    node[right, font=\sffamily\tiny, text=green!60!black] {N tokens} (licensee.north -| sigtok);
  \draw[guidearrow, dashed, orange!60!black] (valtok.south) --
    node[right, font=\sffamily\tiny, text=orange!60!black] {1 token} (licensee.north -| valtok);

  % === Arrows: Licensee deposits to validator ===
  \draw[guidearrow, thick, red!60!black] (licensee.south) -- node[right, font=\sffamily\tiny, text=red!60!black, align=left] {deposit\\sig + val} (wp.north -| licensee);

  % === Note: Ref Token stays with authority ===
  \draw[guidearrow, dotted, gray] (licnft.west) -- ++(-0.8, 0) |- (auth.west)
    node[pos=0.25, left, font=\sffamily\tiny, text=gray!60!black, align=right] {Ref Token (100)\\stays with\\authority};

\end{tikzpicture}
\caption{Token taxonomy showing all four token types, their holders, and their flow through the system. The authority mints all tokens. The licensee receives the User Token (222), signature tokens, and validity tokens. Signing deposits sig+val tokens into the SignatureCollectionValidator. The Reference Token (100) remains with the authority for status updates.}
\label{fig:ch5-tokens}
\end{figure}


% ============================================================================
\section{Wallet Roles}
\label{sec:ch5-wallets}

The system defines four wallet roles, each with distinct permissions and key storage requirements. This role separation is not merely organizational---it is enforced on-chain by minting policies (which check for the authority's key) and by validator scripts (which check for specific signer public key hashes). A wallet's role determines what transactions it can construct and sign.

\subsection{Authority}

The authority wallet belongs to the licensing board, professional association, or credentialing body. It is the most privileged role in the system. The authority can \textbf{mint and burn all token types} (license NFTs, signature tokens, validity tokens), \textbf{create minting policies} (the native scripts that restrict minting to the authority's key), and \textbf{deploy validator scripts} (the Plutus V2 scripts that manage signature collection and dues enforcement). The authority also controls the Reference Token (label 100) of every CIP-68 license NFT it has issued, giving it unilateral revocation power.

Because the authority key can mint tokens that represent professional credentials, its compromise would be catastrophic: an attacker could issue fraudulent licenses. For this reason, authority keys are stored \textbf{server-side with AES-256-GCM encryption}, protected by restrictive file permissions, and ideally backed by hardware security modules in production. The migration path to multi-signature governance (3-of-5 \texttt{ScriptNofK} or KERI AID delegation) distributes this trust across multiple key holders, but the single-key model is the current implementation.

\subsection{Licensee}

The licensee is the professional---the engineer, physician, attorney, or other credentialed individual. The licensee \textbf{holds the license NFT User Token (222)} in their personal wallet, \textbf{receives and deposits signature and validity tokens} during signing ceremonies, and \textbf{pays dues} to the authority's dues enforcement contract.

Critically, the licensee wallet uses a \textbf{CIP-30 non-custodial model}. The system never has access to the licensee's private keys. Instead, the licensee connects a standard Cardano browser wallet (Eternl, Nami, or Lace) to the system's web interface. When a transaction requires the licensee's signature---for example, depositing tokens into a validator---the application constructs the transaction and sends it to the CIP-30 wallet for signing. The licensee reviews and approves the transaction in their wallet extension. This is the same pattern used by all Cardano DeFi applications and means that the system cannot spend the licensee's tokens without their explicit consent.

\subsection{Signer}

The signer role is functionally similar to the licensee role but is used specifically in the context of multi-party signing ceremonies. A signer \textbf{deposits signature and validity tokens} into the \texttt{SignatureCollectionValidator} to attest to a work product. Like the licensee, the signer uses a \textbf{CIP-30 non-custodial wallet} and retains full control of their private keys.

In many cases, the licensee and signer are the same entity---a PE signing their own drawings is both licensee (they hold the license) and signer (they deposit tokens into the validator). The role distinction matters primarily in multi-party scenarios where different individuals contribute signatures to the same work product, each with their own license NFT and token balances.

\subsection{Observer}

The observer is any third party that needs to verify a credential but does not need to create, hold, or transfer tokens. Observers have \textbf{read-only access}: they can query the blockchain (via Blockfrost or any other chain indexer) to check whether a license NFT exists, whether its Reference Token datum shows \texttt{status: active}, and whether the policy ID matches a registered authority. Observers need only a \textbf{public key}---no private key, no wallet software, no account.

The observer role is what makes the system publicly verifiable. A hospital verifying a physician's license, a building department checking an engineer's credentials, or a client confirming an attorney's bar membership all act as observers. The verification is permissionless: no registration, no API key, no bilateral agreement with the issuing authority is required.

\begin{figure}[htbp]
\centering
\begin{tikzpicture}[
  guidebox/.style={draw, rounded corners=4pt, thick, align=center, font=\sffamily\small, fill=white, minimum height=1cm},
  guidearrow/.style={-{Latex[length=2.5mm]}, thick, color=gray!70!black},
  guidelabel/.style={font=\sffamily\footnotesize, text=gray!50!black},
]
  % === Four Roles ===
  \node[guidebox, fill=red!6, minimum width=4cm, minimum height=2.4cm, text width=3.6cm] (authority) at (-4, 4)
    {\textbf{Authority}\\[3pt]
    \footnotesize Mint/burn all tokens\\Create minting policies\\Deploy validators\\Update CIP-68 datums\\Finalize work products\\[3pt]
    \scriptsize \textit{Key: AES-256-GCM\\server-side}};

  \node[guidebox, fill=blue!8, minimum width=4cm, minimum height=2.4cm, text width=3.6cm] (licensee) at (4, 4)
    {\textbf{Licensee}\\[3pt]
    \footnotesize Hold License NFT (222)\\Deposit sig/val tokens\\Pay annual dues\\Request token refills\\[3pt]
    \scriptsize \textit{Key: CIP-30 wallet\\(Eternl/Nami/Lace)}};

  \node[guidebox, fill=green!8, minimum width=4cm, minimum height=2.4cm, text width=3.6cm] (signer) at (-4, -0.5)
    {\textbf{Signer}\\[3pt]
    \footnotesize Deposit sig+val tokens\\into validators\\Attest to work products\\Multi-party ceremonies\\[3pt]
    \scriptsize \textit{Key: CIP-30 wallet\\(Eternl/Nami/Lace)}};

  \node[guidebox, fill=yellow!8, minimum width=4cm, minimum height=2.4cm, text width=3.6cm] (observer) at (4, -0.5)
    {\textbf{Observer}\\[3pt]
    \footnotesize Read-only verification\\Query policy IDs\\Check Reference Token\\datum (CIP-31)\\[3pt]
    \scriptsize \textit{Key: Public key only\\No wallet needed}};

  % === Central: Blockchain ===
  \node[guidebox, fill=purple!8, minimum width=3cm, minimum height=1cm, text width=2.6cm] (chain) at (0, 1.8)
    {\textbf{Cardano}\\Blockchain};

  % === Arrows ===
  \draw[guidearrow, red!60!black] (authority.south east) -- node[above, font=\sffamily\tiny, sloped] {mint / revoke} (chain.north west);
  \draw[guidearrow, blue!60!black] (licensee.south west) -- node[above, font=\sffamily\tiny, sloped] {hold / deposit} (chain.north east);
  \draw[guidearrow, green!60!black] (signer.north east) -- node[below, font=\sffamily\tiny, sloped] {deposit sig+val} (chain.south west);
  \draw[guidearrow, dashed, yellow!60!black] (observer.north west) -- node[below, font=\sffamily\tiny, sloped] {read-only verify} (chain.south east);

  % === Inter-role arrows ===
  \draw[guidearrow, dotted, gray] (authority.east) -- node[above, font=\sffamily\tiny] {issues tokens to} (licensee.west);
  \draw[guidearrow, dotted, gray] (licensee.south) -- ++(0,-0.5) -| node[near end, right, font=\sffamily\tiny, align=left] {same entity\\or different} (signer.east);

\end{tikzpicture}
\caption{Four wallet roles and their interactions with the Cardano blockchain. The authority has full mint/burn privileges with server-side key storage. Licensees and signers use non-custodial CIP-30 wallets. Observers perform read-only verification with no wallet required.}
\label{fig:ch5-roles}
\end{figure}


% ============================================================================
\section{Credential Lifecycle}
\label{sec:ch5-lifecycle}

A credential moves through eight distinct steps from creation to verification (and potentially revocation). Each step involves specific participants, specific transactions, and specific on-chain effects. This section traces the full journey in detail.

\subsection{Step (a): Policy Creation}

Before any tokens can be minted, the authority creates a \textbf{minting policy}. In this system, the minting policy is a \textbf{native script}---specifically, a \texttt{ScriptPubkey} that requires the authority's payment verification key hash. The policy script is serialized, and its Blake2b-224 hash becomes the \textbf{Policy ID}, a 28-byte identifier that uniquely and deterministically identifies all tokens minted under this policy.

The Policy ID is critical for verification. When a third party sees a license NFT in someone's wallet, the first thing they check is the Policy ID. If the Policy ID matches one registered to a known licensing board, the token is legitimate (it could only have been minted by someone who controls the authority's private key). If the Policy ID is unknown, the token is either fraudulent or from an unrecognized authority. This is analogous to checking the seal on a physical certificate---except that forging a Policy ID requires forging an Ed25519 private key, which is computationally infeasible.

The policy creation step is a one-time operation per authority. The resulting Policy ID is published, registered, and used for all subsequent minting operations. No on-chain transaction is required to create the policy itself---the policy exists as a serialized script that is included in minting transactions.

\subsection{Step (b): License NFT Minting (CIP-68)}

The authority constructs a minting transaction that creates two tokens simultaneously under the same Policy ID:

\begin{enumerate}[nosep]
  \item The \textbf{Reference Token (label 100)} is sent to the authority's own UTxO. Its inline datum contains the credential metadata:
  \begin{itemize}[nosep]
    \item \texttt{license\_type}: ``PE'', ``MD'', ``JD'', etc.
    \item \texttt{jurisdiction}: ``US-CA'', ``US-NY'', etc.
    \item \texttt{issue\_date}: ISO 8601 date string
    \item \texttt{licensee\_identity}: public key hash of the licensee
    \item \texttt{status}: ``active'' (initial value)
  \end{itemize}
  \item The \textbf{User Token (label 222)} is sent to the licensee's wallet address. This token has the same Policy ID and base asset name as the Reference Token, differing only in the CIP-68 label prefix.
\end{enumerate}

The transaction includes the native script as a witness and is signed by the authority's payment signing key. Upon confirmation, the licensee's wallet shows the User Token, and the authority's address holds the Reference Token with its inline datum. The total cost is approximately 2.2 ADA (0.2 ADA transaction fee + 2.0 ADA minimum UTxO deposit for the Reference Token).

\subsection{Step (c): Signature and Validity Token Minting}

In a separate transaction (or batched into the same transaction for efficiency), the authority mints \textbf{signature tokens} and a \textbf{validity token} for the licensee. Signature tokens are minted in a batch---typically 100 at a time---under the same Policy ID. The validity token is minted with an asset name or datum encoding its expiry slot. Both are sent to the licensee's wallet address. The cost is approximately 2.2 ADA per batch.

\subsection{Step (d): Token Holding}

The licensee now holds three types of tokens in their CIP-30 wallet: the User Token (222), a balance of signature tokens, and a validity token. These tokens are visible in any Cardano wallet that supports native assets. The licensee can see their token balances, verify the Policy ID, and confirm that the tokens are genuine. No further action is required until the licensee needs to sign a document or renew their validity.

\subsection{Step (e): Signing (Token Deposit)}

When the licensee (acting as a signer) needs to sign a document, they deposit one signature token and one validity token into the \texttt{SignatureCollectionValidator} script address. The transaction includes a \texttt{SignerDatum} containing:

\begin{itemize}[nosep]
  \item The signer's public key hash (PKH)
  \item The SHA-256 hash of the document being signed
  \item The Policy ID of the signer's license
  \item The validity token's expiry slot
  \item The current slot (deposit timestamp)
\end{itemize}

The validator script checks that: (1) the signature token has the correct Policy ID, (2) the validity token has the correct Policy ID and its expiry slot is in the future, and (3) the transaction is signed by the signer whose PKH is in the datum. If all checks pass, the transaction is accepted, and the tokens are locked at the script address. Each signer creates an \textbf{independent UTxO} at the script address, avoiding the eUTxO contention problem that would arise if all signers tried to spend the same UTxO.

\subsection{Step (f): Verification}

Any observer can verify the signing event by querying the validator script address and checking the UTxOs. The verification process (detailed in Section~\ref{sec:ch5-verification}) confirms that tokens with the correct Policy ID are present at the validator address, that the Policy ID maps to a registered authority, that the Reference Token datum shows \texttt{status: active}, and that the validity token had not expired at the time of deposit. This verification can be performed by anyone, at any time, without permission from any party.

\subsection{Step (g): Renewal}

When the validity token approaches expiry, the licensee pays dues to the authority's \textbf{dues enforcement contract}. This Plutus V2 contract verifies that the correct ADA amount was paid and that the payment is within the grace period. Upon successful payment, the authority mints a new validity token with an updated expiry slot and sends it to the licensee's wallet. The old validity token is either burned or simply becomes useless (its expiry slot is in the past, so no validator will accept it).

The renewal process creates a direct economic link between professional standing and signing capacity. If a licensee fails to pay dues, their validity token expires, and they cannot sign documents---even though they still hold the User Token (222) and may have remaining signature tokens. This mirrors real-world licensing: a PE whose license lapses due to unpaid dues cannot legally stamp drawings, even though they once held a valid license.

\subsection{Step (h): Revocation}

If the authority needs to revoke a credential (due to disciplinary action, fraud, or other cause), it updates the Reference Token's inline datum to \texttt{status: revoked}. This is a simple transaction: the authority spends the UTxO holding the Reference Token and creates a new UTxO at the same address with the same token but an updated datum. The transaction costs approximately 0.3 ADA and requires only the authority's signature.

The revocation is effective immediately upon confirmation (within approximately 20 seconds for the first block, with practical finality in approximately 5 minutes). Any subsequent verification will read the Reference Token's datum and find \texttt{status: revoked}. The User Token (222) remains in the licensee's wallet, but it is now meaningless---a verified credential requires both the User Token \emph{and} an active Reference Token datum. This mechanism is detailed in Section~\ref{sec:ch5-revocation}.

\begin{figure}[p]
\centering
\begin{tikzpicture}[
  guidebox/.style={draw, rounded corners=4pt, thick, align=center, font=\sffamily\small, fill=white, minimum height=1cm},
  guidearrow/.style={-{Latex[length=2.5mm]}, thick, color=gray!70!black},
  guidelabel/.style={font=\sffamily\footnotesize, text=gray!50!black},
  stepbox/.style={draw, rounded corners=4pt, thick, align=left, font=\sffamily\scriptsize, fill=white, minimum height=1.4cm, text width=4.2cm, inner sep=5pt},
]
  % Title
  \node[font=\sffamily\large\bfseries, text=gray!70!black] at (0, 15.5) {Credential Lifecycle --- 8 Steps};

  % === Step (a) ===
  \node[stepbox, fill=gray!10] (stepa) at (-3.5, 14)
    {\textbf{(a) Policy Creation}\\[2pt]
    Authority creates native script\\
    Blake2b-224 hash $\to$ Policy ID\\
    One-time per authority};

  % === Step (b) ===
  \node[stepbox, fill=blue!8] (stepb) at (3.5, 14)
    {\textbf{(b) Mint License NFT}\\[2pt]
    CIP-68 dual token pair:\\
    Ref Token (100) $\to$ authority\\
    User Token (222) $\to$ licensee\\
    Cost: $\sim$2.2 ADA};

  % === Step (c) ===
  \node[stepbox, fill=green!8] (stepc) at (-3.5, 11.5)
    {\textbf{(c) Mint Sig + Val Tokens}\\[2pt]
    Signature tokens (qty N)\\
    Validity token (slot expiry)\\
    Both sent to licensee\\
    Cost: $\sim$2.2 ADA each};

  % === Step (d) ===
  \node[stepbox, fill=yellow!8] (stepd) at (3.5, 11.5)
    {\textbf{(d) Token Holding}\\[2pt]
    Licensee CIP-30 wallet holds:\\
    -- User Token (222)\\
    -- N signature tokens\\
    -- 1 validity token};

  % === Step (e) ===
  \node[stepbox, fill=orange!8] (stepe) at (-3.5, 9)
    {\textbf{(e) Signing (Deposit)}\\[2pt]
    Signer deposits to validator:\\
    1 sig token + 1 val token\\
    + SignerDatum (PKH, doc hash,\\
    \ \ policy ID, expiry, slot)};

  % === Step (f) ===
  \node[stepbox, fill=purple!8, text=white] (stepf) at (3.5, 9)
    {\textbf{(f) Verification}\\[2pt]
    Observer queries validator:\\
    Check Policy ID $\to$ authority\\
    Check Ref Token datum\\
    status == active, not expired};

  % === Step (g) ===
  \node[stepbox, fill=green!8] (stepg) at (-3.5, 6.5)
    {\textbf{(g) Renewal}\\[2pt]
    Licensee pays dues via\\
    dues enforcement contract\\
    Authority mints new validity\\
    token with updated expiry};

  % === Step (h) ===
  \node[stepbox, fill=red!6] (steph) at (3.5, 6.5)
    {\textbf{(h) Revocation}\\[2pt]
    Authority updates Ref Token\\
    datum: status $\to$ ``revoked''\\
    Cost: $\sim$0.3 ADA\\
    Effective in $\sim$20s (1 block)};

  % === Arrows ===
  \draw[guidearrow] (stepa.east) -- (stepb.west);
  \draw[guidearrow] (stepb.south) -- ++(0, -0.3) -| (stepc.north);
  \draw[guidearrow] (stepc.east) -- (stepd.west);
  \draw[guidearrow] (stepd.south) -- ++(0, -0.3) -| (stepe.north);
  \draw[guidearrow] (stepe.east) -- (stepf.west);
  \draw[guidearrow] (stepf.south) -- ++(0, -0.3) -| (stepg.north);
  \draw[guidearrow, dashed, red!60!black] (stepg.east) -- node[above, font=\sffamily\tiny, text=red!50!black] {or} (steph.west);

  % === Participants sidebar ===
  \node[font=\sffamily\scriptsize\bfseries, text=gray!50!black, rotate=90, anchor=south] at (-6.5, 10.5) {PARTICIPANTS};
  \draw[thick, gray!30] (-6.2, 6) -- (-6.2, 15);

  \node[font=\sffamily\tiny, text=red!60!black, anchor=west] at (-6.1, 14.3) {Authority};
  \node[font=\sffamily\tiny, text=red!60!black, anchor=west] at (-6.1, 13.7) {Authority};
  \node[font=\sffamily\tiny, text=red!60!black, anchor=west] at (-6.1, 11.8) {Authority};
  \node[font=\sffamily\tiny, text=blue!60!black, anchor=west] at (-6.1, 11.2) {Licensee};
  \node[font=\sffamily\tiny, text=blue!60!black, anchor=west] at (-6.1, 9.3) {Signer};
  \node[font=\sffamily\tiny, text=gray!60!black, anchor=west] at (-6.1, 8.7) {Observer};
  \node[font=\sffamily\tiny, text=blue!60!black, anchor=west] at (-6.1, 6.8) {Licensee};
  \node[font=\sffamily\tiny, text=red!60!black, anchor=west] at (-6.1, 6.2) {Authority};

  % === Legend ===
  \node[draw, rounded corners=3pt, fill=white, inner sep=4pt, font=\sffamily\tiny, text width=4cm, align=left]
    at (0, 4.5) {
    \textbf{Legend}\\[2pt]
    Solid arrow: normal flow\\
    Dashed red arrow: alternative path\\
    Steps (a)--(d): setup phase\\
    Steps (e)--(f): operational phase\\
    Steps (g)--(h): maintenance phase
    };

\end{tikzpicture}
\caption{Complete credential lifecycle showing all eight steps from policy creation through revocation. The setup phase (a--d) is performed once per credential. The operational phase (e--f) repeats for each signing event. The maintenance phase (g--h) occurs periodically or as needed.}
\label{fig:ch5-lifecycle}
\end{figure}


% ============================================================================
\section{Revocation Mechanism}
\label{sec:ch5-revocation}

This is the most important architectural decision in the entire system. Revocation---the ability for an authority to invalidate a previously issued credential---is the feature that separates a useful credential system from a toy. If a physician loses their license, if an engineer commits fraud, if an attorney is disbarred, the credential must become invalid \emph{immediately and globally}. This section explains why revocation is hard on Cardano, why the naive approach fails, and how the CIP-68 pattern solves the problem.

\subsection{The Problem: eUTxO Ownership}

In Cardano's extended UTxO model, every token sits inside a UTxO, and every UTxO is controlled by either an address (requiring the owner's private key to spend) or a script (requiring satisfaction of the script's conditions to spend). This is a feature, not a bug---it is what makes eUTxO secure. But it creates a specific problem for revocation:

\textbf{You cannot spend (or burn) someone else's UTxO without their private key.}

This is an absolute, non-negotiable property of the ledger. The authority cannot construct a transaction that spends the licensee's UTxO because the ledger requires the licensee's signature for that UTxO. The authority does not have the licensee's private key (and should never have it, by the non-custodial design). This means the authority \emph{cannot forcibly remove a token from the licensee's wallet}.

\subsection{The Broken Approach: Direct Token Burn}

The naive approach to revocation is: ``when the authority wants to revoke a license, burn the licensee's NFT.'' This approach fails completely in eUTxO. To burn a token, you must spend the UTxO that contains it. To spend that UTxO, you need the private key that controls the address where it sits. The licensee's token sits at the licensee's address. The authority does not have the licensee's key. Therefore, the authority cannot burn the licensee's token. Period.

One might suggest requiring the licensee to send the token to the authority upon revocation, but this requires the licensee's cooperation---which is precisely what you cannot count on when revoking someone's license for cause. A physician whose license is being revoked for malpractice is unlikely to voluntarily return their credential token.

Another approach might be to hold the token at a script address controlled by the authority rather than at the licensee's personal address. But this defeats the purpose of the credential: the licensee should be able to prove they hold the token, which requires the token to be in their wallet. If the token is at a script address, the licensee does not ``hold'' it in any meaningful sense.

\subsection{The CIP-68 Solution: Datum-Based Revocation}

The CIP-68 standard resolves this tension elegantly by splitting the credential into two tokens with different custody:

\begin{enumerate}[nosep]
  \item The \textbf{User Token (222)} goes to the licensee's wallet. The authority \emph{cannot} touch it. The licensee holds it forever (or until they voluntarily send it somewhere). This is fine because \textbf{the User Token is not the source of truth}.
  \item The \textbf{Reference Token (100)} stays at the authority's UTxO. The authority \emph{can} spend this UTxO at any time because the authority controls the private key for that address. The Reference Token carries an \textbf{inline datum} with a \texttt{status} field.
\end{enumerate}

Revocation is performed by the authority spending the UTxO that holds the Reference Token and creating a new UTxO at the same address with the same Reference Token but an updated datum: \texttt{status: ``revoked''}. This transaction:

\begin{itemize}[nosep]
  \item Requires \textbf{only the authority's signature}---no interaction with the licensee.
  \item Costs approximately \textbf{0.3 ADA} (a simple transaction fee; no Plutus execution needed for a native script policy).
  \item Takes effect within approximately \textbf{20 seconds} (one block confirmation).
  \item Is \textbf{globally visible}---any observer who reads the Reference Token's inline datum will see \texttt{status: revoked}.
\end{itemize}

The User Token (222) in the licensee's wallet is now meaningless. It still exists---the licensee still ``holds'' it---but any verifier performing a proper verification will check the Reference Token's datum and find \texttt{status: revoked}. The User Token without an active Reference Token is like a physical certificate after the issuing board has published a revocation notice: the paper exists, but the credential it represents has been invalidated.

\subsection{The minUTxO Economics}

Every UTxO on Cardano must carry a minimum ADA deposit (approximately 2.0 ADA for a UTxO carrying a native token with an inline datum). This ADA is not ``spent''---it is locked and can be recovered when the UTxO is consumed. When the authority mints the Reference Token, approximately 2.0 ADA is locked at the UTxO. When the authority performs a revocation (spending and re-creating the UTxO with updated datum), the 2.0 ADA remains locked at the new UTxO. If the authority later decides to \textbf{burn the Reference Token entirely} (removing it from the ledger), the locked ADA is returned to the authority's wallet, minus the transaction fee of approximately 0.3 ADA. The net recovery is approximately 1.7 ADA per burned Reference Token.

For a licensing board managing 10,000 licenses, the total locked ADA at Reference Token UTxOs is approximately 20,000 ADA (approximately \$6,000 at \$0.30/ADA). This is the economic cost of maintaining revocable credentials---a modest amount compared to the annual operating budgets of professional licensing boards.

\begin{figure}[htbp]
\centering
\begin{tikzpicture}[
  guidebox/.style={draw, rounded corners=4pt, thick, align=center, font=\sffamily\small, fill=white, minimum height=1cm},
  guidearrow/.style={-{Latex[length=2.5mm]}, thick, color=gray!70!black},
  guidelabel/.style={font=\sffamily\footnotesize, text=gray!50!black},
]
  % === BEFORE REVOCATION ===
  \node[font=\sffamily\bfseries, text=gray!60!black] at (-4.5, 7.5) {BEFORE Revocation};

  % Authority side - before
  \node[guidebox, fill=gray!10, minimum width=4.2cm, minimum height=1.8cm, text width=3.8cm] (auth_before) at (-6.5, 5.5)
    {\textbf{Authority UTxO}\\[2pt]
    \footnotesize Ref Token (100)\\
    Inline Datum:\\
    \texttt{status: "active"}\\
    \texttt{type: "PE"}\\
    \texttt{jurisdiction: "US-CA"}\\[1pt]
    \scriptsize Locked: $\sim$2.0 ADA};

  % Licensee side - before
  \node[guidebox, fill=blue!8, minimum width=4.2cm, minimum height=1.4cm, text width=3.8cm] (lic_before) at (-6.5, 2.5)
    {\textbf{Licensee Wallet}\\[2pt]
    \footnotesize User Token (222)\\
    + Signature Tokens\\
    + Validity Token};

  % Verification result - before
  \node[guidebox, fill=green!8, minimum width=3cm, text width=2.6cm] (verify_before) at (-2.3, 4)
    {\footnotesize Verification:\\[2pt] \textbf{VALID}};

  \draw[guidearrow, dashed, green!60!black] (auth_before.east) -- (verify_before.north west);
  \draw[guidearrow, dashed, green!60!black] (lic_before.east) -- (verify_before.south west);

  % === Divider ===
  \draw[thick, gray!30, dashed] (0, 1.5) -- (0, 8);

  % === AFTER REVOCATION ===
  \node[font=\sffamily\bfseries, text=red!60!black] at (4.5, 7.5) {AFTER Revocation};

  % Authority side - after
  \node[guidebox, fill=red!6, minimum width=4.2cm, minimum height=1.8cm, text width=3.8cm] (auth_after) at (2.5, 5.5)
    {\textbf{Authority UTxO}\\[2pt]
    \footnotesize Ref Token (100)\\
    Inline Datum:\\
    \texttt{status: "revoked"}\\
    \texttt{type: "PE"}\\
    \texttt{jurisdiction: "US-CA"}\\[1pt]
    \scriptsize Locked: $\sim$2.0 ADA};

  % Licensee side - after (unchanged!)
  \node[guidebox, fill=blue!8, minimum width=4.2cm, minimum height=1.4cm, text width=3.8cm] (lic_after) at (2.5, 2.5)
    {\textbf{Licensee Wallet}\\[2pt]
    \footnotesize User Token (222)\\
    \scriptsize (still present, but now meaningless)};

  % Verification result - after
  \node[guidebox, fill=red!6, minimum width=3cm, text width=2.6cm] (verify_after) at (6.7, 4)
    {\footnotesize Verification:\\[2pt] \textbf{REVOKED}};

  \draw[guidearrow, dashed, red!60!black] (auth_after.east) -- (verify_after.north west);
  \draw[guidearrow, dashed, red!60!black] (lic_after.east) -- (verify_after.south west);

  % === Revocation arrow ===
  \draw[guidearrow, very thick, red!70!black, -{Latex[length=4mm]}]
    (auth_before.south east) .. controls +(1.5, -1) and +(-1.5, -1) ..
    node[below, font=\sffamily\scriptsize, text=red!60!black, align=center]
    {Authority updates\\datum: $\sim$0.3 ADA\\No licensee interaction} (auth_after.south west);

  % === Note ===
  \node[draw, rounded corners=3pt, fill=yellow!8, inner sep=4pt, font=\sffamily\tiny, text width=8cm, align=left]
    at (0, 0.5) {
    \textbf{Key insight:} The licensee's wallet is \emph{never touched}. The authority only modifies
    its own UTxO (the Reference Token). Verifiers always check the Reference Token datum,
    so the User Token alone is insufficient proof of valid licensure.
    };

\end{tikzpicture}
\caption{Before and after revocation. The authority updates only its own Reference Token datum from ``active'' to ``revoked.'' The licensee's User Token remains in their wallet but is now meaningless for verification purposes. No interaction with the licensee is required.}
\label{fig:ch5-revocation}
\end{figure}


% ============================================================================
\section{Verification Flow}
\label{sec:ch5-verification}

Verification is the operation that makes the entire system worthwhile. Everything else---minting, signing, renewal, revocation---exists to produce a state that can be verified. This section describes the exact sequence of steps a third-party observer follows to verify a credential, and explains why each step is necessary.

\subsection{Step-by-Step Verification}

The verification flow consists of seven steps, each building on the previous one. The entire process completes in under 5 seconds, dominated by Blockfrost API latency (200--500ms per query).

\textbf{Step 1: Identify the License NFT.} The verifier obtains the licensee's wallet address (or the licensee presents their wallet via CIP-30). The verifier queries the address's UTxO set and identifies the License NFT User Token (label 222). This step confirms that the licensee holds a token that \emph{claims} to be a license. It does not yet prove anything about the token's legitimacy.

\textbf{Step 2: Extract the Policy ID.} Every Cardano native token has a Policy ID---the Blake2b-224 hash of the minting policy script. The verifier extracts the Policy ID from the User Token. This is a deterministic operation; the Policy ID is part of the token's on-chain identity.

\textbf{Step 3: Verify the Policy ID against a registry.} The verifier checks whether the extracted Policy ID matches any registered authority. This is the step that distinguishes a legitimate credential from a fraudulent one. The registry could be a published list maintained by a standards body, a smart contract on-chain, or a well-known database. If the Policy ID is not recognized, verification fails---the token was minted by an unknown party and cannot be trusted.

\textbf{Step 4: Find the CIP-68 Reference Token.} The verifier constructs the Reference Token's asset name by replacing the User Token's label prefix (222) with the Reference Token's label prefix (100), keeping the same base asset name. The verifier then searches for this Reference Token on-chain. Under CIP-68, the Reference Token should be at a known address (the authority's address or a designated script address).

\textbf{Step 5: Read the inline datum (CIP-31 Reference Input).} Using CIP-31 reference inputs, the verifier reads the Reference Token's inline datum \emph{without spending the UTxO}. This is a read-only operation that does not require any private key and does not create a transaction. The datum contains the credential metadata, including the \texttt{status} field, license type, jurisdiction, issue date, and licensee identity.

\textbf{Step 6: Check status and expiration.} The verifier confirms two conditions: (1) \texttt{status == "active"} (if \texttt{status == "revoked"}, verification fails immediately), and (2) \texttt{expiration > current\_slot} (if the credential has expired, verification fails). The current slot is obtained from the Cardano chain tip.

\textbf{Step 7: Verification complete.} If all checks pass, the credential is verified. The verifier now has cryptographic proof that: the credential was issued by a recognized authority (Policy ID match), the credential is currently active (datum status check), the credential has not expired (slot comparison), and the licensee holds the corresponding User Token (UTxO query).

\begin{figure}[htbp]
\centering
\begin{tikzpicture}[
  guidebox/.style={draw, rounded corners=4pt, thick, align=center, font=\sffamily\small, fill=white, minimum height=1cm},
  guidearrow/.style={-{Latex[length=2.5mm]}, thick, color=gray!70!black},
  guidelabel/.style={font=\sffamily\footnotesize, text=gray!50!black},
  vstep/.style={draw, rounded corners=3pt, thick, fill=white, font=\sffamily\scriptsize, minimum height=0.9cm, minimum width=7cm, align=left, inner sep=5pt},
]
  % === Observer start ===
  \node[guidebox, fill=yellow!8, minimum width=3cm] (observer) at (0, 11.5)
    {\textbf{Observer / Verifier}};

  % === Steps ===
  \node[vstep, fill=blue!8] (s1) at (0, 10)
    {\textbf{1.} Identify License NFT (User Token 222) in holder's wallet};

  \node[vstep, fill=blue!8] (s2) at (0, 8.8)
    {\textbf{2.} Extract Policy ID (Blake2b-224 hash of minting policy)};

  \node[vstep, fill=orange!8] (s3) at (0, 7.6)
    {\textbf{3.} Verify Policy ID matches a registered authority};

  \node[vstep, fill=green!8] (s4) at (0, 6.4)
    {\textbf{4.} Find CIP-68 Reference Token (same policy, label 100)};

  \node[vstep, fill=green!8] (s5) at (0, 5.2)
    {\textbf{5.} Read inline datum via CIP-31 Reference Input (read-only)};

  \node[vstep, fill=purple!8, text=white] (s6) at (0, 4.0)
    {\textbf{6.} Check: status == ``active'' AND expiration > current\_slot};

  \node[vstep, fill=green!8] (s7) at (0, 2.8)
    {\textbf{7.} Verification complete --- credential is valid};

  % === Arrows ===
  \draw[guidearrow] (observer.south) -- (s1.north);
  \draw[guidearrow] (s1.south) -- (s2.north);
  \draw[guidearrow] (s2.south) -- (s3.north);
  \draw[guidearrow] (s3.south) -- (s4.north);
  \draw[guidearrow] (s4.south) -- (s5.north);
  \draw[guidearrow] (s5.south) -- (s6.north);
  \draw[guidearrow] (s6.south) -- (s7.north);

  % === Failure paths ===
  \node[guidebox, fill=red!6, minimum width=2cm, text width=1.8cm, font=\sffamily\scriptsize] (fail1) at (5.5, 7.6)
    {FAIL:\\Unknown\\Policy ID};
  \draw[guidearrow, red!60!black] (s3.east) -- node[above, font=\sffamily\tiny, text=red!50!black] {no match} (fail1.west);

  \node[guidebox, fill=red!6, minimum width=2cm, text width=1.8cm, font=\sffamily\scriptsize] (fail2) at (5.5, 4.0)
    {FAIL:\\Revoked or\\Expired};
  \draw[guidearrow, red!60!black] (s6.east) -- node[above, font=\sffamily\tiny, text=red!50!black] {status $\neq$ active} (fail2.west);

  % === Timing ===
  \node[draw, rounded corners=3pt, fill=yellow!8, inner sep=4pt, font=\sffamily\tiny, text width=6cm, align=left]
    at (0, 1.5) {
    \textbf{Timing:} Total $<$ 5 seconds\\
    Steps 1--2: $<$ 500ms (UTxO query)\\
    Step 3: $<$ 100ms (local registry lookup)\\
    Steps 4--5: $<$ 500ms (asset query + datum read)\\
    Step 6: $<$ 100ms (comparison)
    };

\end{tikzpicture}
\caption{Seven-step verification flow. An observer can verify a credential in under 5 seconds with no permission required. Failure at Step 3 (unknown Policy ID) or Step 6 (revoked/expired status) immediately terminates verification with a negative result.}
\label{fig:ch5-verification}
\end{figure}

\subsection{Permissionless Verification}

A crucial property of this verification flow is that it is \textbf{permissionless}. The observer does not need:

\begin{itemize}[nosep]
  \item An account or API key with the issuing authority
  \item A bilateral trust agreement with any party
  \item Any special software beyond a Blockfrost API client (or any Cardano chain indexer)
  \item The licensee's cooperation (the on-chain state is public)
\end{itemize}

This is fundamentally different from current verification infrastructure, where a hospital must contact the specific state medical board, provide credentials, wait for a response, and trust that the board's database is accurate and current. In the blockchain-based system, the verification is mathematical: the Policy ID either matches or it does not, the datum either shows ``active'' or it does not, and the expiration either exceeds the current slot or it does not.

\subsection{Verification Without a Full Node}

The verification flow described above uses the Blockfrost API, which means the verifier trusts Blockfrost to return accurate chain data. For higher-assurance verification, the observer can run their own Cardano node and query the ledger directly, eliminating all third-party trust. For an intermediate assurance level, Mithril snapshot proofs (planned for 2026) will allow light clients to verify chain state with cryptographic proofs without running a full node.


% ============================================================================
\section{Cost Economics}
\label{sec:ch5-costs}

Understanding the on-chain costs is essential for evaluating the system's economic viability. Every Cardano transaction has two cost components: the \textbf{transaction fee} (paid to stake pool operators, approximately 0.17--0.30 ADA depending on transaction size and Plutus execution) and the \textbf{minimum UTxO deposit} (ADA locked at the output UTxO, recoverable when the UTxO is consumed, approximately 1.5--2.0 ADA depending on the datum and token payload). The table below summarizes costs for each operation.

\begin{table}[htbp]
\centering
\caption{On-chain cost breakdown for credential operations}
\label{tab:ch5-costs}
\small
\begin{tabular}{@{}p{4.5cm}ccc@{}}
\toprule
\textbf{Operation} & \textbf{Tx Fee (ADA)} & \textbf{minUTxO (ADA)} & \textbf{Total (ADA)} \\
\midrule
Mint License NFT (CIP-68 pair) & $\sim$0.2 & $\sim$2.0 & $\sim$2.2 \\
Mint Signature Tokens (batch of 100) & $\sim$0.2 & $\sim$2.0 & $\sim$2.2 \\
Mint Validity Token & $\sim$0.2 & $\sim$2.0 & $\sim$2.2 \\
Sign Document (deposit to validator) & $\sim$0.2 & $\sim$2.0 & $\sim$2.2 \\
Finalize Work Product (Plutus V2) & $\sim$0.3 & $\sim$2.0 & $\sim$2.3 \\
Revoke License (datum update) & $\sim$0.3 & 0\textsuperscript{$\dagger$} & $\sim$0.3 \\
Renew Validity (new token mint) & $\sim$0.2 & $\sim$2.0 & $\sim$2.2 + dues \\
Burn Reference Token (cleanup) & $\sim$0.3 & $-$2.0\textsuperscript{$\ddagger$} & net $\sim$$-$1.7 \\
\bottomrule
\end{tabular}

\vspace{4pt}
\footnotesize
\textsuperscript{$\dagger$}Revocation modifies the existing Reference Token UTxO in place; the $\sim$2.0 ADA deposit remains locked.\\
\textsuperscript{$\ddagger$}Burning the Reference Token releases the locked minUTxO, yielding a net refund of $\sim$1.7 ADA after fees.
\end{table}

\subsection{Comparison to Ethereum}

The same operations on Ethereum would cost significantly more due to the ERC-721 contract execution model. Minting an NFT on Ethereum requires deploying or interacting with a smart contract, with gas costs that vary dramatically based on network congestion. During periods of high activity, gas prices have exceeded 500 gwei, pushing a single NFT mint above \$100. Even at moderate gas prices, a typical ERC-721 mint costs \$1--15, and complex operations like multi-party signature collection (requiring multiple contract calls and state updates) can cost \$5--50 per operation. Cardano's deterministic fee model eliminates this variability entirely: the cost of a transaction is known before submission and does not change based on network load.

Solana offers lower absolute costs (sub-\$0.01 per transaction) but with a different risk profile: multiple network outages, a proof-of-history consensus mechanism with less formal verification than Ouroboros, and a programming model (Rust-based programs) less suited to formal verification of credential logic.

\subsection{Annual Cost Projection}

For a deployment managing 10,000 professional licenses with standard operational assumptions (one license mint per licensee, 100 signature tokens per licensee per year, one validity renewal per year, 10 signing operations per licensee per year), the annual on-chain costs at current ADA prices (\$0.30/ADA) are:

\begin{table}[htbp]
\centering
\caption{Annual cost projection for 10,000 licenses at \$0.30/ADA}
\label{tab:ch5-annual}
\small
\begin{tabular}{@{}p{5cm}cc@{}}
\toprule
\textbf{Category} & \textbf{ADA} & \textbf{USD} \\
\midrule
Initial license minting (one-time, year 1) & 22,000 & \$6,600 \\
Signature token minting (annual batch) & 22,000 & \$6,600 \\
Validity token renewal (annual) & 22,000 & \$6,600 \\
Signing operations (100,000/year) & 220,000 & \$66,000 \\
Revocations (est. 2\% = 200/year) & 60 & \$18 \\
\midrule
\textbf{Year 1 Total} & \textbf{286,060} & \textbf{\$85,818} \\
\textbf{Recurring Annual} & \textbf{264,060} & \textbf{\$79,218} \\
\bottomrule
\end{tabular}

\vspace{4pt}
\footnotesize
Note: The bulk of the cost is signing operations. If each licensee performs only 1 signing per year (as in annual license renewal), the recurring cost drops to approximately \$6,300--8,400 annually. Verification queries are free (read-only Blockfrost API calls). Hydra state-channel batching (2026 roadmap) could reduce signing costs by 90\%+.
\end{table}

The key economic insight is that \textbf{verification is free}. Reading on-chain state through Blockfrost or any chain indexer does not require a transaction and incurs no ADA cost. The current verification infrastructure charges \$5--25 per query (FSMB, Nursys, NCEES verification services). For an organization performing thousands of verifications annually, the savings from free verification alone can offset the entire on-chain transaction cost.

\subsection{Locked ADA and Capital Efficiency}

The minUTxO requirement means that ADA is locked at token-carrying UTxOs. For 10,000 licenses, approximately 20,000 ADA ($\sim$\$6,000) is locked at Reference Token UTxOs, plus additional ADA locked at signature and validity token UTxOs held by licensees. This locked ADA is not lost---it can be recovered by consuming the UTxOs---but it represents an opportunity cost. CIP-68 datum patterns and potential future Hydra state channels could reduce locked ADA requirements by over 90\% through token bundling and off-chain settlement.


% ============================================================================
\section{Chapter Summary}
\label{sec:ch5-summary}

This chapter has presented the complete credential verification system built on Cardano. The key takeaways are:

\begin{enumerate}[leftmargin=*]
  \item \textbf{Three-layer architecture:} The blockchain provides settlement finality, the application layer provides business logic, and the SQLite database provides fast queries. Blockfrost bridges the application to the chain.

  \item \textbf{Four token types:} License NFTs (CIP-68 dual-token pair), signature tokens (fungible, consumed during signing), validity tokens (time-bounded, linked to dues), and work products (logical composites coordinated by validators).

  \item \textbf{Four wallet roles:} Authority (full minting privileges, server-side keys), licensee (non-custodial CIP-30 wallet), signer (non-custodial CIP-30 wallet), and observer (read-only, public key only).

  \item \textbf{Eight lifecycle steps:} Policy creation, license minting, token minting, token holding, signing, verification, renewal, and revocation---each with specific participants, costs, and on-chain effects.

  \item \textbf{CIP-68 revocation:} The authority controls the Reference Token (100) and can update its datum to \texttt{status: revoked} unilaterally, without touching the licensee's wallet. This is the core architectural innovation that makes on-chain revocation possible in eUTxO.

  \item \textbf{Permissionless verification:} Any observer can verify a credential in under 5 seconds by checking the Policy ID against a registry and reading the Reference Token's inline datum. No account, no API key, no bilateral agreement required.

  \item \textbf{Economically viable:} Annual on-chain costs for 10,000 licenses range from \$6,300 (minimal signing) to \$85,000 (heavy signing), with verification queries free. This compares favorably to current verification service fees of \$5--25 per query.
\end{enumerate}

The next chapter examines the smart contract layer in detail, including the Plutus V2 validator logic, redeemer patterns, and the security properties that make the system resistant to the six attack surfaces identified in the threat model.
